% Generated draft: requires mathematical and source review.
\documentclass[11pt]{article}
\usepackage[margin=25mm]{geometry}
\usepackage{fontspec}
\setmainfont{Latin Modern Roman}
\usepackage{amsmath,amssymb,mathtools,unicode-math}
\setmathfont{Latin Modern Math}
\usepackage{enumitem,needspace,xurl,hyperref}
\setlength{\parindent}{0pt}
\setlength{\parskip}{5pt}
\setlength{\emergencystretch}{3em}
\newcommand{\N}{\mathbb{N}}
\newcommand{\Z}{\mathbb{Z}}
\newcommand{\Q}{\mathbb{Q}}
\newcommand{\R}{\mathbb{R}}
\newcommand{\C}{\mathbb{C}}
\newcommand{\F}{\mathbb{F}}
\newcommand{\A}{\mathbb{A}}
\newcommand{\PP}{\mathbb{P}}
\newcommand{\E}{\mathbb{E}}
\newcommand{\eps}{\varepsilon}
\newcommand{\dd}{\,\mathrm{d}}
\begin{document}
\section*{UnsolvedMath ID: 700610}
\textbf{Source problem code:} catalogue entry
\subsection{Definitions and mathematical statement}
For a group \(G\) and elements \(x,y\in G\), let \(\langle x,y\rangle\) denote the subgroup they generate. Interpret the source term \emph{binary finite} as
\[
\forall x,y\in G,\qquad |\langle x,y\rangle|<\infty.
\]
The supplied record does not define this term, so this interpretation requires confirmation against the cited source. A group \(G\) is \emph{finitely generated} if \(G=\langle S\rangle\) for some finite subset \(S\subseteq G\). It is \emph{residually finite} if, for every \(g\in G\setminus\{1\}\), there are a finite group \(F\) and a homomorphism \(\varphi:G\to F\) such that \(\varphi(g)\ne 1\).

For a prime \(p\), a \(p\)-subgroup means a subgroup each of whose elements has order \(p^k\) for some integer \(k\geq 0\). A \emph{Sylow \(p\)-subgroup} is understood here to mean a \(p\)-subgroup maximal under inclusion among the \(p\)-subgroups of \(G\). The supplied record does not specify its infinite-group Sylow convention; this convention also requires source confirmation.

Under these terminology conventions, the question attributed to V. P. Shunkov is: does there exist a group \(G\) such that \(G\) is infinite, finitely generated, residually finite, and binary finite, and
\[
\forall p\text{ prime}\;\forall P\leq G,\qquad
P\text{ a Sylow }p\text{-subgroup of }G
\ \Longrightarrow\ |P|<\infty?
\]

The source additionally reports that A. V. Rozhkov proved the corresponding existence assertion when every Sylow subgroup is required only to be \emph{locally finite}, meaning that each of its finitely generated subgroups is finite. This is a supplied historical assertion, attributed to a 1997 doctoral dissertation, and is not independently verified here.
\subsection{Short English statement}
Does an infinite finitely generated residually finite group exist in which every two elements generate a finite subgroup and every Sylow subgroup is finite? The interpretations of ``binary finite'' and ``Sylow subgroup'' stated above require source confirmation. The record reports that Rozhkov obtained existence with locally finite Sylow subgroups instead.
\subsection{Sources}
\begin{itemize}
\item Local UnsolvedMath registry, permanent ID \(700610\), exact source problem code \texttt{catalogue entry}, title \emph{Kourovka Notebook Problem 9.77}; supplied record revision \texttt{d371ee917ac9647fbb06cd6fc529deec6f3dda82}, updated \texttt{2026-10-16T14:00:00Z}. The registry labels the problem \emph{open}; that label is supplied curation, not an independently verified mathematical assessment.
\item \emph{The Kourovka Notebook (Unsolved Problems in Group Theory)}, original problem 9.77, ninth issue (1984), proposed by V. P. Shunkov. The supplied background identifies the twenty-first issue edition (2026), arXiv:1401.0300v48. The supplied collection description names E. I. Khukhro and V. D. Mazurov as editors.
\item A. V. Rozhkov, doctoral dissertation (1997), cited in the supplied statement for the existence result with locally finite Sylow subgroups. No dissertation title, institution, page reference, or link is supplied.
\end{itemize}
\end{document}
